Los resultados obtenidos durante la implementación del entorno experimental confirman la pertinencia del uso de plataformas basadas en honeypots, como T-Pot, para el análisis y la detección de intrusiones dentro de redes locales aisladas, especialmente cuando el objetivo se orienta a la observación detallada del comportamiento de los atacantes más que a la prevención inmediata, ya que la progresión de los ataques simulados, desde fases iniciales de reconocimiento hasta intentos de explotación activa y denegación de servicio, evidenció patrones coherentes con modelos ampliamente documentados en la literatura, como la Cyber Kill Chain, lo cual refuerza tanto la validez metodológica del escenario planteado como la utilidad del enfoque de observación pasiva adoptado en este estudio.\\
De manera complementaria, la capacidad de T-Pot para capturar, correlacionar y visualizar eventos en tiempo casi real se consolidó como un elemento clave para la comprensión integral de las tácticas, técnicas y procedimientos empleados por los atacantes, dado que la integración de múltiples honeypots con herramientas como Suricata, Kibana y utilidades de monitoreo de recursos permitió no solo registrar los intentos de intrusión, sino también analizar su impacto sobre el sistema y clasificar distintos perfiles de amenaza en función de su comportamiento, aspecto particularmente relevante en contextos académicos y organizacionales donde las amenazas internas suelen pasar desapercibidas y, en consecuencia, pueden generar efectos significativos si no se detectan de manera oportuna.\\
En términos generales, los hallazgos confirman que un entorno controlado y éticamente diseñado constituye un marco eficaz para evaluar tecnologías de detección de intrusiones sin comprometer activos reales ni exponer infraestructuras productivas, aun cuando se reconoce que el aislamiento del laboratorio limita la extrapolación directa de los resultados a escenarios con tráfico externo real y ataques no controlados; sin embargo, estas limitaciones no invalidan los aportes del estudio, sino que abren la posibilidad de trabajos futuros orientados a entornos híbridos y a la integración con sistemas de prevención, por lo que, en conclusión, la implementación de T-Pot en una red LAN aislada demostró ser una estrategia efectiva para la detección, el análisis y la clasificación de intrusiones simuladas, contribuyendo al fortalecimiento de estrategias proactivas de ciberseguridad y resaltando el valor de los honeypots como mecanismos de defensa basados en el entendimiento profundo del comportamiento del atacante.